The rapid growth of open source software (OSS) has significantly shaped modern software development, traditionally emphasizing tools, libraries, and infrastructure that benefit developer communities. However, a growing subset of OSS projects is shifting focus toward broader societal impact, where developers leverage their technical expertise to address real-world challenges such as healthcare accessibility, disaster response, and social inclusion. These projects, referred to as Open Source Software for Social Good (OSS4SG), represent an emerging paradigm in which software development is not only a technical endeavor but also a means of contributing to societal well-being. Despite their increasing presence, OSS4SG projects remain relatively underexplored, particularly in understanding how they differ from conventional OSS in terms of goals, contributors, and impact.

This paper seeks to bridge that gap by examining the motivations, decision-making processes, and challenges faced by contributors to OSS4SG. Unlike traditional OSS contributors, who are often driven by skill development or career advancement, OSS4SG contributors are more strongly motivated by the desire to create meaningful societal impact and address pressing social issues. At the same time, these projects introduce unique complexities, including difficulties in project discovery, coordination across diverse stakeholders, and sustaining long-term engagement. By systematically analyzing these aspects, the study provides insights into how OSS4SG differs from general OSS and highlights opportunities to better support contributors and organizations involved in socially impactful software development.

\section{Motivation}
The motivation for this paper stems from a fundamental gap in how open source software (OSS) is traditionally understood versus how it is increasingly being used in practice. While OSS has long been associated with building developer-centric tools and infrastructure, there is a growing movement where developers intentionally contribute to projects aimed at addressing societal challenges such as healthcare access, disaster response, and social inclusion. Despite the clear societal relevance of these efforts, there is little structured understanding of how such projects—termed Open Source Software for Social Good (OSS4SG)—differ from conventional OSS in terms of contributor motivations, project characteristics, and impact. The authors recognize that existing OSS research primarily focuses on technical or self-oriented incentives (e.g., skill development, reputation building), leaving a gap in understanding contributions driven by social impact.

A key motivation of the study is to investigate whether contributing to OSS4SG is fundamentally different from contributing to traditional OSS. Prior work in fields such as social work and education suggests that initiatives tied to societal impact can attract diverse participants and foster a sense of urgency and purpose. However, it remains unclear whether these dynamics translate into the open source ecosystem. The paper is therefore motivated to explore questions such as: Do contributors to OSS4SG have different motivations? How do they identify and select such projects? And what unique challenges do they encounter? By addressing these questions, the authors aim to uncover whether OSS4SG represents merely a subset of OSS or a distinct category with its own behavioral and social dynamics.

Another important motivation is the lack of a clear, contributor-informed definition of “social good” within the context of open source. Although the term has been used across disciplines—including sociology, AI, and software engineering—it remains loosely defined and inconsistently applied in OSS communities. This ambiguity makes it difficult for contributors to identify relevant projects and for researchers to systematically study them. The paper is therefore motivated to formalize the concept of OSS4SG by grounding it in contributors’ perspectives, particularly focusing on projects that explicitly target societal issues and benefit specific communities. Establishing such a definition is critical for enabling further research and improving discoverability of these projects.

Finally, the study is driven by practical challenges observed in the OSS4SG ecosystem. Contributors often struggle to discover impactful projects, evaluate their credibility, and sustain long-term engagement due to factors such as limited funding, unclear project direction, and coordination across diverse stakeholders. At the same time, organizations and project maintainers lack insights into how to attract and retain contributors who are motivated by social impact. By systematically analyzing motivations, selection criteria, and barriers, the paper aims to generate actionable insights that can improve the design, visibility, and sustainability of OSS4SG projects. Ultimately, the motivation is not only to understand this emerging space but also to support its growth as a meaningful avenue for leveraging software engineering for societal benefit.

\section{Methodology}
The paper adopts an exploratory sequential mixed-methods approach to systematically investigate contributors’ experiences in Open Source Software for Social Good (OSS4SG). The study begins with qualitative exploration and is followed by quantitative validation, allowing the authors to both discover themes and generalize findings at scale. Specifically, the methodology is structured into two main phases: semi-structured interviews and a large-scale survey, ensuring both depth and breadth in understanding OSS4SG contributors .

\paragraph{Phase 1: Semi-Structured Interviews}
The qualitative phase focuses on understanding contributors’ lived experiences. The authors first identify OSS4SG projects using curated platforms such as Ovio and Digital Public Goods, which classify projects based on social impact domains like healthcare, education, and sustainability. From these sources, they recruit \textbf{21 participants across 437 OSS4SG projects}, ensuring diversity by sampling across different levels of project popularity and activity. Interviews are conducted remotely, lasting 45–55 minutes, and follow a semi-structured format to allow flexibility in probing participant experiences. Key topics include \textbf{contributors’ backgrounds, motivations, perceptions of OSS4SG vs OSS, project selection criteria, and challenges}. All interviews are recorded, transcribed, and analyzed using \textbf{inductive thematic analysis}, where recurring patterns are identified through iterative coding and refinement using tools like ATLAS.ti. This phase enables the authors to derive core themes that later inform the survey design .

\paragraph{Data Analysis (Qualitative)}
The analysis follows a rigorous multi-step process. First, researchers independently generate \textbf{open codes from interview transcripts}, capturing recurring ideas and observations. These codes are then collaboratively refined into a structured codebook, ensuring consistency and reliability. Next, relationships between codes are identified and grouped into higher-level themes such as motivations, definitions of social good, and challenges in OSS4SG. The process continues iteratively until \textbf{theoretical saturation} is reached, meaning no new themes emerge from additional data. This systematic approach ensures that findings are grounded in participant experiences and not preconceived assumptions .

\paragraph{Phase 2: Survey Study}
To validate and generalize insights from the interviews, the authors design a \textbf{28-question survey} targeting both OSS4SG contributors and general OSS contributors. The survey is distributed to contributors from two groups: (1) OSS4SG projects (437 curated projects) and (2) general OSS projects randomly sampled from GitHub. A total of \textbf{5765 invitations are sent, resulting in 517 valid responses}, which is consistent with typical response rates in software engineering research. The survey captures data on \textbf{demographics, motivations, project selection factors, perceptions of OSS4SG, and challenges}. It also includes structured Likert-scale questions and scenario-based questions (e.g., impact trade-offs) to quantitatively assess contributor preferences .

\paragraph{Data Analysis (Quantitative)}
The survey data is analyzed using statistical techniques to compare OSS4SG and OSS contributors. Respondents are categorized into groups based on their perceptions and experience (e.g., those who distinguish OSS4SG from OSS vs those who do not). The authors compute \textbf{descriptive statistics and perform hypothesis testing (e.g., Wilcoxon rank-sum tests)} to identify significant differences in motivations and behaviors. Additionally, a \textbf{logistic regression model} is used to examine how demographic factors influence perceptions of OSS4SG. This quantitative phase validates qualitative findings and highlights statistically significant differences between OSS4SG and traditional OSS contributors .

\textbf{Research Questions Framing}
The entire methodology is guided by four central research questions:
\begin{enumerate}[label=\textbf{RQ\arabic*},leftmargin=*]
    \item How do contributors define OSS4SG?
    \item What motivates contributors to participate in OSS4SG?
    \item How do contributors select OSS4SG projects?
    \item What challenges do contributors face in OSS4SG?
\end{enumerate}
These questions structure both the interview protocol and survey design, ensuring alignment between qualitative exploration and quantitative validation.

\section{Key Findings}
The study reveals that Open Source Software for Social Good (OSS4SG) is not merely a subset of traditional OSS, but a \textbf{distinct category with unique characteristics in definition, motivation, project selection, and challenges}. A majority of contributors (over 52\%) explicitly distinguish OSS4SG from general OSS, primarily based on whether a project \textbf{targets societal issues and benefits specific communities}. Contributors emphasize that OSS4SG is defined not just by openness, but by its \textbf{intentional social impact and identifiable beneficiaries}, rather than indirect or downstream effects .

A central finding is the difference in contributor motivations. While traditional OSS contributors are often driven by \textbf{self-oriented goals such as skill development, portfolio building, or career advancement}, OSS4SG contributors are significantly more motivated by \textbf{altruistic and impact-driven goals}. These include helping underserved populations, solving real-world problems, and achieving emotional fulfillment by “leaving a fingerprint on the world.” Quantitative results reinforce this distinction, showing that OSS4SG contributors place \textbf{significantly less importance on learning new skills or improving career prospects} and \textbf{significantly more importance on solving societal issues}.

When selecting projects, contributors to OSS4SG evaluate factors differently compared to traditional OSS. While common factors like project activity and clarity still matter, OSS4SG contributors place a much stronger emphasis on \textbf{trust in project owners, alignment with personal values, and the societal relevance of the project}. They are also less influenced by traditional signals such as project popularity. Additionally, contributors tend to prioritize projects that \textbf{have global reach, long-term impact, and benefit people they personally relate to or know}. This highlights a more value-driven and impact-sensitive decision-making process.

The study also uncovers several challenges unique to OSS4SG. The most prominent issue is the \textbf{difficulty in discovering and matching contributors with relevant OSS4SG projects}, due to the lack of standardized labels or discovery mechanisms. Contributors also face \textbf{funding instability}, which affects sustained participation, and \textbf{coordination challenges} arising from diverse stakeholders, including non-technical participants such as NGOs and government bodies. Additionally, differences in technical expertise and cultural backgrounds can lead to \textbf{communication barriers and conflicts}, further complicating collaboration.

\section{Implications}

The findings lead to several actionable implications for improving the OSS4SG ecosystem. First, there is a clear need to improve project discoverability. The authors suggest introducing \textbf{standardized badges, labels, or classification systems} within platforms like GitHub to explicitly identify OSS4SG projects. This would help bridge the gap between contributors who are motivated to help and projects that need support. Additionally, curated directories and nomination systems should adopt clearer criteria based on \textbf{target users and societal impact}, making it easier to identify meaningful projects.

Second, project maintainers and organizations should better communicate the purpose and impact of their projects. Since OSS4SG contributors are highly motivated by social impact, it is crucial to \textbf{explicitly highlight the societal goals, intended beneficiaries, and expected outcomes} in project documentation. Providing transparent information about project ownership and organizational backing can also build \textbf{trust}, which is a key factor in contributor decision-making. Emphasizing these aspects can significantly improve contributor engagement and retention.

Another important implication relates to sustainability and support structures. OSS4SG projects often struggle with funding and long-term commitment, so there is a need for \textbf{more stable funding models and institutional support}. Organizations, non-profits, and governments can play a critical role in sustaining these projects by providing resources, visibility, and infrastructure. Furthermore, creating clearer onboarding processes and contribution guidelines can help reduce entry barriers for new contributors.

Finally, the study highlights the need for stronger safeguards in OSS4SG projects due to their sensitive nature. Many projects deal with vulnerable populations or critical societal issues, which introduces risks related to \textbf{privacy, ethics, and user safety}. The authors recommend implementing \textbf{training, moderation mechanisms, and clear reporting structures} to handle conflicts and protect both contributors and end users. Additionally, projects involving sensitive data or collaboration with public institutions must ensure \textbf{compliance with regulations and ethical standards}.